% !TITLE 涉江采芙蓉
% !AUTHOR 无名氏
% !DYNASTY 两汉
% !GENRE 五言古诗
% !ORDER 4

\begin{poem}
涉江采芙蓉\textsuperscript{1}，兰泽\textsuperscript{2}多芳草。
采之欲遗\textsuperscript{3}谁？所思在远道\textsuperscript{4}。
还顾\textsuperscript{5}望旧乡，长路漫浩浩\textsuperscript{6}。
同心而离居，忧伤以终老。
\end{poem}

\begin{annotations}
\notetext{1}{芙蓉：荷花的别称。古人采芙蓉常赠送给思念的亲人或配偶，象征纯洁美好的情感。}
\notetext{2}{兰泽：生长着兰草的水泽。兰，一种香草。}
\notetext{3}{遗：遗（wèi），赠送。}
\notetext{4}{远道：遥远的地方。}
\notetext{5}{还顾：回头看。}
\notetext{6}{漫浩浩：路途极其遥远、漫无边际的样子。}
\end{annotations}

\begin{appreciation}
芙蓉，即荷花，古人以芙蓉为高洁的象征。如离骚里所写“制芰荷以为衣兮，集芙蓉以为裳”。兰花同样是楚辞里常见的香草意象。而“长路漫浩浩”，即“长路漫漫浩浩”，屈原有“路漫漫其修远兮”的句子，“漫浩浩”极言路途之遥远。这句诗也有着楚辞里的痕迹。因此我们不难想到这首诗的作者很有可能是楚地的文人，他继承了楚文化的文学语言，又融入了汉乐府朴实平易，写出了这样一首意蕴悠长、让人回味无穷的小诗。

关于这首诗，历来有个众说纷纭的疑惑，这位“涉江采芙蓉”的人，到底是远方的游子还是在家的思妇？

有人认为这首诗运用了对写法。采芙蓉的人是在家的思妇，她想象远方的游子此时此刻应该在思念家乡，思念着在家的自己，正如自己思念远方的游子一般，此为“同心”。两人相隔万里，此为“离居”。两人“同心而离居”，自然因相思而悲伤了，可谓“忧伤以终老”。

也有人认为采芙蓉的是远方的游子。因为采香草这样充满象征意味的行为在当时应当是一位有一定社会地位与文化的“士”才会做的事情。这位漂泊士人在长满兰草的水泽里望着家乡的方向，可他穷极目力也望不到遥远的家乡。于是他感叹“同心而离居，忧伤以终老”。

不过或许这首诗的主人公到底是谁也并不那么重要。虽然他们的姓名已经消失在历史的长河里了，但他们的相思依旧寄托在那岁岁枯荣的荷花上、幽香芬芳的兰泽里。
\end{appreciation}
